\begin{pastexam}[label={px:nccu-108-dsa-06}]{政大資科所 108 學年度 第 6 題}
（問答）

(10\%) Give an application of the heap data structure.
\end{pastexam}

\begin{pastsoln}
題目只要一個應用，寫一個並說清楚「堆積在裡面做什麼、為什麼比別的結構好」。

\textbf{應用：優先佇列}（資結書第~8.3~節）。 堆積最主要的用途是實作優先佇列：一直有新的工作加進來，
每次要取出優先權最高（鍵最小或最大）的那一個。 二元堆積插入、取出最小值都是 $O(\lg n)$，看最小值
$O(1)$。 對照：用沒排序的陣列，插入 $O(1)$ 但取出要掃一遍 $\Theta(n)$；用排好序的陣列，取出 $O(1)$ 但
插入要搬移 $\Theta(n)$。 兩種操作交錯很多次時，堆積的 $O(\lg n)$ 好得多。

具體的例子是\textbf{作業系統的工作排程}：每個工作有一個優先權，新的工作隨時送進來（\proc{Insert}），
CPU 空下來就取出優先權最高的工作來執行（\proc{Extract-Max}）。 $n$ 個工作在佇列裡時，每次排程只要
$O(\lg n)$。

其他常見的用法（任選一個寫也可以）：
\begin{itemize}
\item \textbf{堆積排序}（資結書第~14.6~節）：先在 $O(n)$ 時間建堆積，再取出 $n$ 次最大值，$O(n \lg n)$，
而且是原地排序。
\item \textbf{Dijkstra 最短路徑、Prim 最小生成樹}（資結書第~25.2、24.3~節）：每一步取出距離（或連接成本）
最小的頂點，用二元堆積總時間 $O((|V| + |E|) \lg |V|)$。
\item \textbf{Huffman 編碼}（資結書第~19.3~節）：每次取出兩個頻率最小的節點合併，再放回去。
\item \textbf{找最大的 $k$ 個數}：維持一個大小 $k$ 的最小堆積，看到比堆頂大的數就取代堆頂，$O(n \lg k)$。
\item \textbf{$k$ 路合併}（外部排序）：$k$ 個排好的串列，堆積裡放每個串列目前的開頭，每次取出最小的，
$O(n \lg k)$。
\end{itemize}
\end{pastsoln}
